
\subsection{Java annotation-ийн үндсэн ойлголт}

Annotation нь Java хэл дээр класс, интерфэйс, гишүүн функц, гишүүн өгөгдөл, параметр зэрэг программын элементүүдэд нэмэлт мэдээлэл буюу metadata өгөх зориулалттай хэлний механизм юм. Annotation нь өөрөө бизнес логик хэрэгжүүлдэггүй боловч компилятор, хөгжүүлэлтийн хэрэгсэл, framework болон runtime орчинд тухайн элементийн зориулалт, төлөв, хэрэглээний талаар мэдээлэл өгдөг.\footnote{Oracle. (2024).
\textit{Lesson: Annotations}.
The Java\texttrademark{} Tutorials.
\url{https://docs.oracle.com/javase/tutorial/java/annotations/}}

Java annotation-уудыг зориулалт, хэрэглээнээс нь хамааруулан built-in annotation, meta-annotation, custom annotation гэж ангилан авч үздэг.

\subsubsection{Built-in annotation}

Built-in annotation нь Java хэл өөрөө агуулдаг annotation-ууд бөгөөд ихэвчлэн компилятор болон хөгжүүлэгчид зориулсан байдаг. Эдгээр нь framework-ийн түвшний annotation биш, харин Java compiler-д нэмэлт заавар, тайлбар өгөх зорилготой.

Түгээмэл built-in annotation-уудад дараах annotation-ууд орно:
\begin{itemize}
    \item \texttt{@Override}
    \item \texttt{@Deprecated}
    \item \texttt{@SuppressWarnings}
    \item \texttt{@FunctionalInterface}
\end{itemize}

Жишээлбэл, \texttt{@Override} annotation нь эцэг классаас өвлөсөн функцыг дахин тодорхойлж байгааг илэрхийлдэг бол \texttt{@Deprecated} annotation нь тухайн элемент хуучирсан, цаашид хэрэглэхээс зайлсхийх шаардлагатайг тэмдэглэдэг.

\subsubsection{Meta-annotation}

Meta-annotation нь өөр annotation тодорхойлоход ашиглагддаг annotation юм. Өөрөөр хэлбэл annotation дээр хэрэглэгддэг annotation гэж үзэж болно. Java хэлэнд custom annotation боловсруулах үед meta-annotation-ууд чухал үүрэгтэй байдаг. Энэхүү дипломын ажлын хүрээнд хэрэгжүүлэх framework-д custom annotation үүсгэхдээ эдгээр meta-annotation-уудыг ашиглах шаардлагатай.

Түгээмэл meta-annotation-уудыг Хүснэгт \ref{tab:meta-annotation} -д үзүүлэв.

\begin{table}[H]
\centering
\small
\setlength{\tabcolsep}{4pt}
\caption{Java хэлний түгээмэл meta-annotation-ууд тэдгээрийн үүрэг}
\label{tab:meta-annotation}
\begin{tabularx}{\textwidth}{|L{0.24\textwidth}|Y|}
\hline
\textbf{Annotation} & \textbf{Үүрэг} \\
\hline
\texttt{@Target} & Annotation-ийг ямар төрлийн элемент дээр ашиглаж болохыг тодорхойлно \\
\hline
\texttt{@Retention} & Annotation ямар хугацаанд хадгалагдахыг тодорхойлно \\
\hline
\texttt{@Documented} & Annotation нь JavaDoc баримт бичигт харагдах эсэхийг заана \\
\hline
\texttt{@Inherited} & Annotation нь удамшлын үед subclass руу дамжих эсэхийг тодорхойлно \\
\hline
\texttt{@Repeatable} & Нэг annotation-ийг олон удаа хэрэглэх боломж олгоно \\
\hline
\end{tabularx}
\end{table}

\newpage
\subsubsection{Custom annotation}

Custom annotation нь meta-annotation ашиглан хэрэглэгч өөрийн шаардлагад нийцүүлэн тодорхойлсон annotation юм. Framework боловсруулах үед component, service, repository, inject зэрэг зориулалтаар custom annotation үүсгэх нь түгээмэл байдаг. Ийм annotation-ууд нь runtime үед reflection ашиглан уншигдаж, тухайн класст ямар үйлдэл хийхийг тодорхойлоход ашиглагдана.\footnote{Oracle. (2024).
\textit{Annotations} (Java SE 17 \& JDK 17).
\url{https://docs.oracle.com/en/java/javase/17/docs/api/
java.base/java/lang/annotation/Annotation.html}}

Жишээлбэл, дараах \texttt{@Component} annotation нь class түвшинд хэрэглэгдэх бөгөөд runtime үед хадгалагдахаар тодорхойлогдсон байна.

\begin{lstlisting}[language=Java, caption={Custom annotation жишээ}, label={lst:custom-annotation}]
@Target(ElementType.TYPE)
@Retention(RetentionPolicy.RUNTIME)
public @interface Component {
}
\end{lstlisting}




\subsection{Annotation-ий бүтцийн ангилал}

Бүтцийн хувьд annotation-ийг дараах үндсэн төрлүүдэд ангилж үздэг.

\begin{enumerate}
    \item \textbf{Marker annotation} -- parameter агуулаагүй annotation.
    \item \textbf{Single-value annotation} -- нэг утга хүлээн авдаг annotation.
    \item \textbf{Full annotation} -- олон attribute буюу олон утга тодорхойлох боломжтой annotation.
\end{enumerate}

\newpage
\subsubsection{Marker annotation}

Marker annotation нь ямар нэг parameter агуулаагүй, зөвхөн тухайн элементийн төлөв эсвэл зориулалтыг тэмдэглэх зориулалттай байдаг.

\begin{lstlisting}[language=Java, caption={Marker annotation жишээ}, label={lst:marker-annotation}]
@Override
public String toString() {
    return "Example";
}
\end{lstlisting}

\subsubsection{Single-value annotation}

Single-value annotation нь нэг үндсэн параметр авдаг.

\begin{lstlisting}[language=Java, caption={Single-value annotation жишээ}, label={lst:single-value-annotation}]
@SuppressWarnings("unchecked")
public void process() {
}
\end{lstlisting}

\subsubsection{Full annotation}

Full annotation нь олон attribute агуулж болдог тул илүү дэлгэрэнгүй тохиргоо өгөх боломжтой.

\begin{lstlisting}[language=Java, caption={Full annotation жишээ}, label={lst:full-annotation}]
@RequestMapping(value="/users", method=RequestMethod.GET)
public List<User> getUsers() {
    return users;
}
\end{lstlisting}

\subsection{Custom annotation үүсгэхэд ашиглагддаг үндсэн annotation-ууд}

\subsubsection{\texttt{@Target}}

\texttt{@Target} annotation нь тухайн annotation-ийг ямар төрлийн элемент дээр хэрэглэхийг тодорхойлдог. Өөрөөр хэлбэл annotation-ийн хэрэглэгдэх хүрээг зааж өгдөг.

\begin{table}[H]
\centering
\small
\setlength{\tabcolsep}{4pt}
\caption{\texttt{@Target} annotation-д ашиглагддаг \texttt{ElementType} болон тайлбарууд}
\label{tab:elementtype}
\begin{tabularx}{\textwidth}{|L{0.26\textwidth}|Y|}
\hline
\textbf{ElementType} & \textbf{Тайлбар} \\
\hline
\texttt{TYPE} & Class, interface, enum зэрэг төрөл дээр \\
\hline
\texttt{FIELD} & Классын талбар дээр \\
\hline
\texttt{METHOD} & Метод дээр \\
\hline
\texttt{PARAMETER} & Методын параметр дээр \\
\hline
\texttt{CONSTRUCTOR} & Constructor дээр \\
\hline
\texttt{LOCAL\_VARIABLE} & Local variable дээр \\
\hline
\texttt{ANNOTATION\_TYPE} & Өөр annotation дээр \\
\hline
\texttt{PACKAGE} & Package дээр \\
\hline
\end{tabularx}
\end{table}

Жишээлбэл, дараах annotation нь зөвхөн class болон interface зэрэг төрөл дээр хэрэглэгдэхээр тодорхойлогдсон байна.\footnote{Oracle. (2024). \textit{Annotation Type Target}.
Java SE 17 API Documentation.
\url{https://docs.oracle.com/en/java/javase/17/docs/api/
java.base/java/lang/annotation/Target.html}}

\begin{lstlisting}[language=Java, caption={\texttt{@Target} ашигласан жишээ}, label={lst:target-example}]
@Target(ElementType.TYPE)
public @interface Component {
}
\end{lstlisting}

\newpage
\subsubsection{\texttt{@Retention}}

\texttt{@Retention} annotation нь тухайн annotation хэзээ хүртэл хадгалагдахыг тодорхойлдог. Энэ нь annotation-ийн амьдралын хугацааг заадаг бөгөөд compile-time болон runtime орчинд annotation ашиглах эсэхийг шийддэг.\footnote{Oracle. (2024). \textit{Annotation Type Target}.
Java SE 17 API Documentation.
\url{https://docs.oracle.com/en/java/javase/17/docs/api/
java.base/java/lang/annotation/Retention.html}}

Retention policy нь дараах 3 төрөлтэй.

\begin{itemize}
    \item \texttt{SOURCE} -- Annotation нь зөвхөн source code түвшинд оршиж, compile хийх үед \texttt{.class} файл руу орохгүй устдаг.
    \item \texttt{CLASS} -- Annotation нь compile хийсний дараа \texttt{.class} файлд хадгалагдах боловч runtime үед reflection-оор уншигдахгүй.
    \item \texttt{RUNTIME} -- Annotation нь runtime үед хадгалагдаж, reflection ашиглан унших боломжтой.
\end{itemize}

Эдгээрийг Хүснэгт \ref{tab:retention-policy} -д харуулав.

\begin{table}[H]
\centering
\caption{Retention policy-ийн орчны ялгаа}
\label{tab:retention-policy}
\begin{tabular}{|p{3cm}|c|c|c|}
\hline
\textbf{Policy} & \textbf{Source code} & \textbf{Class file} & \textbf{Runtime} \\
\hline
\texttt{SOURCE} & \checkmark & \texttimes & \texttimes \\
\hline
\texttt{CLASS} & \checkmark & \checkmark & \texttimes \\
\hline
\texttt{RUNTIME} & \checkmark & \checkmark & \checkmark \\
\hline
\end{tabular}
\end{table}

\newpage
Framework боловсруулах үед хамгийн чухал нь \texttt{RUNTIME} policy юм.
Учир нь runtime үед reflection ашиглан annotation-ийг уншиж,
dependency injection болон component scan зэрэг механизмыг хэрэгжүүлэх
шаардлагатай байдаг.

Жишээ нь:

\begin{lstlisting}[language=Java, caption={\texttt{@Retention} ашигласан жишээ}, label={lst:retention-example}]
@Target(ElementType.TYPE)
@Retention(RetentionPolicy.RUNTIME)
public @interface Component {
}
\end{lstlisting}

\subsubsection{\texttt{@Documented}}

\texttt{@Documented} annotation нь тухайн custom annotation JavaDoc баримт бичигт харагдах эсэхийг тодорхойлдог. Хэрэв annotation нь албан ёсны API баримт бичигт харагдах шаардлагатай бол энэ annotation-ийг хэрэглэнэ.

\begin{lstlisting}[language=Java, caption={\texttt{@Documented} ашигласан жишээ}, label={lst:documented-example}]
@Documented
@Target(ElementType.TYPE)
@Retention(RetentionPolicy.RUNTIME)
public @interface Service {
}
\end{lstlisting}

\subsubsection{\texttt{@Inherited}}

\texttt{@Inherited} annotation нь class түвшний annotation удамшлын дагуу subclass-д дамжих эсэхийг тодорхойлдог. Энэ нь зөвхөн class дээр хэрэглэгдсэн annotation-д үйлчилдэг.

\begin{lstlisting}[language=Java, caption={\texttt{@Inherited} ашигласан жишээ}, label={lst:inherited-example}]
@Inherited
@Target(ElementType.TYPE)
@Retention(RetentionPolicy.RUNTIME)
public @interface Component {
}
\end{lstlisting}

\subsubsection{\texttt{@Repeatable}}

\texttt{@Repeatable} annotation нь нэг төрлийн annotation-ийг нэг элемент дээр олон удаа хэрэглэх боломж олгодог. Энэ нь ижил төрлийн олон metadata хадгалах шаардлагатай үед хэрэгтэй байдаг.

\begin{lstlisting}[language=Java, caption={\texttt{@Repeatable} ашигласан жишээ}, label={lst:repeatable-example}]
@Repeatable(Roles.class)
public @interface Role {
    String value();
}
\end{lstlisting}

\newpage

\subsection{Spring Framework дахь annotation}

Spring Framework нь annotation-д суурилсан хөгжүүлэлтийн загварыг өргөн ашигладаг бөгөөд annotation-уудыг зориулалтаас нь хамааран component annotation, configuration annotation, dependency injection annotation, web annotation, AOP annotation, transaction annotation зэрэг ангиллаар авч үздэг.\footnote{Spring Framework Documentation. (2024).
\textit{Annotation-based Container Configuration}.
\url{https://docs.spring.io/spring-framework/reference/
core/beans/annotation-config.html}}

\subsubsection{Component annotation}

Component annotation нь Spring container-т bean болгох class-ыг тодорхойлоход ашиглагдана.

Жишээ нь:
\begin{itemize}
    \item \texttt{@Component} -- Package scanning үед framework энэ annotation-той class-уудыг илрүүлнэ. Энэхүү annotation-ийг IoC Container-д ``Энэ классыг чи удирдаарай, объектыг нь үүсгээрэй'' гэж хэлэхэд ашиглана. Тухайн annotation байгаа классыг IoC Container-т бүртгэхэд ашиглана.
    \item \texttt{@Service} -- Business logic давхаргын class тэмдэглэх бөгөөд \texttt{@Component}-ийн тусгай хэлбэр юм.
    \item \texttt{@Repository} -- Өгөгдлийн сангийн давхаргын class тодорхойлоход ашиглана.
    \item \texttt{@Controller} -- Application хэсэгт end-point тэмдэглэх, илэрхийлэх annotation юм.
\end{itemize}

\newpage
\subsubsection{Configuration annotation}

Configuration annotation нь container-т тохиргоог үүсгэж тодорхойлж өгөхөд ашиглагддаг.

Жишээ нь:
\begin{itemize}
    \item \texttt{@Configuration} -- Тохиргооны класс тэмдэглэнэ. Bean method-уудыг агуулна. \texttt{@Component}-ийн тусгай хэлбэр учир container-т бүртгэгддэг.
    \item \texttt{@ComponentScan} -- Хаана package scan хийхийг заана. \texttt{value}, \texttt{basePackage} буюу заасан package-с уншина.
    \item \texttt{@Bean} -- Annotation-тай үүсгэх объектуудыг тодорхойлно.
    \item \texttt{@Import} -- Өөр тохиргооны классыг оруулахад ашиглана. Том project-т тохиргоог хэд хэдэн class-т хуваан \texttt{@Import}-оор нэгтгэнэ.
    \item \texttt{@Scope} -- Bean-ийн буюу тухайн annotation байгаа объектын амьдралын мөчлөгийг тодорхойлно.
\end{itemize}

\subsubsection{Dependency Injection annotation}

Dependency Injection annotation нь dependency injection буюу хамаарлыг илэрхийлэхэд ашиглагддаг.

Жишээ нь:
\begin{itemize}
    \item \texttt{@Autowired} -- Field болон constructor алинд inject хийх вэ гэдгийг тодорхойлно. Spring-д ажилладаг.
    \item \texttt{@Inject} -- Dependency injection хийх field эсвэл constructor заана. Ямар ч JSR-330 container-т ажиллана.
    \item \texttt{@Qualifier} -- Нэг interface-ийг хэд хэдэн class implement хийвэл аль нэгийг нь сонгон inject хийхэд ашиглана. Аль нэгийг нь \texttt{@Qualifier("stripeGateway")} гэж сонгож өгөх боломжтой.
    \item \texttt{@Primary} -- \texttt{@Qualifier} байхгүй бол түүний implementation хийх өгөгдлийг энэ annotation-ийг тэмдэглэн илэрхийлж өгдөг.
\end{itemize}

Жишээ:

\begin{lstlisting}[language=Java, caption={\texttt{@Qualifier} ашигласан жишээ}, label={lst:qualifier-example}]
@Component("mysqlRepo")
public class MySQLUserRepository implements UserRepository {}

@Component("mongoRepo")
public class MongoUserRepository implements UserRepository {}

@Autowired
@Qualifier("mysqlRepo") // аль нэгийг нь заана
private UserRepository userRepo;
\end{lstlisting}

\begin{lstlisting}[language=Java, caption={\texttt{@Primary} ашигласан жишээ}, label={lst:primary-example}]
@Primary // @Qualifier байхгүй үед энэ сонгогдоно
public class StripePaymentGateway implements PaymentGateway { ... }

@Service
public class CheckoutService {
    @Autowired
    private PaymentGateway gateway; // -> StripePaymentGateway сонгогдоно
}
\end{lstlisting}

\newpage
\subsubsection{Web annotation}

Web annotation нь веб аппликэйшин хийхэд ашиглагддаг.

Жишээ нь:
\begin{itemize}
    \item \texttt{@RequestMapping} -- HTTP route тодорхойлох үндсэн annotation. Бүх HTTP method-д ажиллана (\texttt{GET}, \texttt{POST}, \texttt{PUT}, \texttt{DELETE} гэх мэт).
    \item \texttt{@GetMapping} -- \texttt{GET} хүсэлтэд тодорхойлох annotation.
    \item \texttt{@PostMapping} -- \texttt{POST} хүсэлтэд тодорхойлох annotation.
    \item \texttt{@DeleteMapping} -- \texttt{DELETE} хүсэлтэд тодорхойлох annotation.
    \item \texttt{@UpdateMapping} -- \texttt{UPDATE} хүсэлтэд тодорхойлох annotation.
    \item \texttt{@RestController} -- \texttt{@Controller} + \texttt{@ResponseBody} хосолсон annotation. Бүх method-ийн буцаах утгыг JSON болгон хөрвүүлнэ.
\end{itemize}

\subsubsection{AOP annotation}

AOP annotation нь огтлолцсон логик буюу cross-cutting concern-уудыг хэрэгжүүлэхэд ашиглагдана.

Жишээ нь:
\begin{itemize}
    \item \texttt{@Aspect}
    \item \texttt{@Before}
    \item \texttt{@After}
    \item \texttt{@Around}
\end{itemize}

\subsubsection{Transaction annotation}

Transaction annotation нь transaction удирдахад ашиглагдана.

Жишээ нь:
\begin{itemize}
    \item \texttt{@Transactional}
\end{itemize}

\newpage
\subsection{Framework-д хэрэгжүүлэх annotation-ууд}

\subsubsection{Дипломд заавал хэрэгжүүлэх annotation-ууд (DI Framework-ийн цөм)}

\begin{itemize}
    \item \texttt{@Component}, \texttt{@Service}, \texttt{@Repository} -- Bean бүртгэлт
    \item \texttt{@Autowired}, \texttt{@Qualifier}, \texttt{@Primary} -- Dependency injection
    \item \texttt{@Scope} -- Bean амьдралын мөчлөг
\end{itemize}

\subsubsection{Хэрэгжүүлэх боломжтой annotation-ууд (нэмэлт)}

\begin{itemize}
    \item \texttt{@GetMapping}, \texttt{@PostMapping} -- Энгийн HTTP routing
    \item \texttt{@Configuration},  \texttt{@ComponentScan}, \texttt{@Bean}, -- Тохиргоо
    \item \texttt{@Transactional} -- Энгийн transaction
    \item \texttt{@Before}, \texttt{@After} -- Энгийн AOP logging
\end{itemize}

\subsection{Java ба Spring annotation-ийн ялгаа}

Java хэл дээр annotation нь программын элементүүдэд metadata нэмэх зориулалттай бөгөөд built-in annotation, meta-annotation болон custom annotation зэрэг төрлүүдэд ангилагддаг. Харин Spring Framework дээр annotation-ийг хэрэглээний зориулалтаар нь component annotation, configuration annotation, dependency injection annotation болон web annotation зэрэг ангиллаар ашигладаг. Энэхүү судалгааны ажилд annotation болон reflection механизмд тулгуурлан бүрэлдэхүүн хэсгүүдийг автоматаар илрүүлэх dependency injection framework-ийг хэрэгжүүлсэн.

\newpage
\subsection{Compile болон Runtime ойлголт}

\textbf{Compile} гэдэг нь Java эх кодыг \texttt{javac} хөрвүүлэгч ашиглан \texttt{.class} файл болгон хөрвүүлэх үйл явц юм. Энэ үед compiler нь syntax алдааг шалгаж, Java Virtual Machine (JVM)-д ойлгогдох bytecode үүсгэдэг.

\textbf{Runtime} гэдэг нь хөрвүүлэгдсэн \texttt{.class} файлыг JVM санах ойд ачаалж, программыг ажиллуулах үе шат юм. Reflection механизм нь runtime үед ажилладаг бөгөөд annotation унших, class-ийн бүтэц шалгах, объект үүсгэх, метод дуудах зэрэг үйлдлүүдийг энэ шатанд гүйцэтгэдэг.

Иймээс annotation-г хэрэгжүүлэхийн тулд annotation-үүдийг \texttt{RetentionPolicy.}\allowbreak\texttt{RUNTIME} тохиргоотой тодорхойлох шаардлагатай байдаг.

\subsection{Reflection механизм}
Reflection нь Java хэлний runtime орчинд class, method, field, constructor зэрэг бүтцийн мэдээллийг шалгах, унших, өөрчлөх боломж олгодог механизм юм. Энэ механизм нь annotation-д суурилсан framework-ийн дотоод ажиллагаанд чухал үүрэгтэй.

Reflection ашигласнаар дараах боломжууд бүрдэнэ:
\begin{itemize}
    \item class-ийн мэдээлэл унших,
    \item field болон method илрүүлэх,
    \item annotation байгаа эсэхийг шалгах,
    \item constructor ашиглан объект үүсгэх,
    \item field-д утга оноох,
    \item method-ийг runtime үед дуудах
\end{itemize}


% \begin{figure}[H]
%     \centering
%     \includegraphics[width=0.8\textwidth]{images/reflection-flow.png}
%     \caption{Reflection ашиглан annotation уншиж dependency холбох ерөнхий үйл явц}
%     \label{fig:reflection-flow}
% \end{figure}
Энэхүү дипломын ажлын хэрэгжүүлэлтэд reflection механизмыг ашиглан custom annotation-уудтай class-уудыг илрүүлэх, объект үүсгэх, dependency-үүдийг автоматаар холбох үндсэн ажиллагааг хэрэгжүүлнэ.

\subsection{Annotation ба Reflection-ийн уялдаа холбоо}

Annotation болон reflection нь framework-ийн дотоод механизмд салшгүй холбоотой ажилладаг. Annotation нь класст metadata өгдөг бол reflection нь runtime үед уг metadata-г уншиж, түүнд үндэслэн шаардлагатай үйлдлийг гүйцэтгэдэг.

Жишээлбэл, хэрэв класс \texttt{@Component} annotation-той бол framework нь reflection ашиглан уг annotation-ийг илрүүлж, тухайн классыг container-д бүртгэж, шаардлагатай үед объект үүсгэх боломжтой болно. Үүнтэй адил \texttt{@Inject} annotation тэмдэглэгдсэн талбарыг reflection ашиглан олж, холбогдох dependency-г автоматаар оноож болно.

Иймээс annotation нь ``юу хийх'' талаар мэдээлэл өгдөг бол reflection нь ``түүнийг хэрхэн хэрэгжүүлэх'' механизмыг хангадаг гэж үзэж болно.